\section{Online Routing Algorithms}
\label{sec:online}

In the online setting, we must make irrevocable routing decisions without knowledge of future arrivals or exact output lengths. We propose a unified framework that combines primal-dual algorithm design with learning-augmented predictions.

\subsection{Primal-Dual Framework}

We treat the subscription capacity as a scarce resource and assign it a dynamic \textit{shadow price}. This price represents the opportunity cost of using a subscription slot.

\textbf{Quota Shadow Price ($\psi$).}
For the daily quota constraint ($S_Q$), let $z = k/Q$ be the fraction of quota consumed. We define an exponential threshold function $\psi(z)$:
\begin{equation}
    \psi(z) = L \cdot \left( \frac{U}{L} \right)^z
\end{equation}
where $L$ and $U$ are estimated lower and upper bounds on the request value (avoided API cost). A request is routed to $S_Q$ only if its value exceeds this threshold.

\begin{theorem}[Competitive Ratio]
The primal-dual algorithm with threshold $\psi(z)$ achieves a competitive ratio of $O(\ln(U/L))$ for the online knapsack problem with known values.
\end{theorem}

The exponential threshold ensures that the opportunity cost of quota grows smoothly from $L$ (when quota is abundant) to $U$ (when quota is nearly exhausted). Early requests with moderate value are accepted; late requests must have high value to justify consuming scarce quota. This logarithmic competitive ratio is known to be optimal for the online knapsack problem.

\textbf{Concurrency Shadow Price ($\lambda$).}
For the concurrency constraint ($S_C$), we define a congestion price $\lambda_t$ based on the current utilization of the concurrency slots. When the system is saturated (queue full or utilization high), $\lambda_t$ increases, admitting only requests with high \textit{value density} (value per unit time).
\begin{equation}
    \lambda_t = \begin{cases}
    0 & \text{if capacity available} \\
    \min_{r \in \text{Queue}} \frac{v(r)}{w(r) \cdot p(r)} & \text{if saturated}
    \end{cases}
\end{equation}

\subsection{Learning-Augmented Decision Making}

A key challenge is that the ``value'' of a request (its output token count $n_i^{\text{out}}$) is unknown at arrival time. Standard estimators (e.g., historical averages) can be risky: over-estimating value leads to wasting scarce quota on cheap requests.

We employ a \textbf{Learning-Augmented} approach with two components. First, we use \textbf{quantile regression} to train a lightweight model that predicts the distribution of output tokens. Specifically, we predict the 10th percentile (P10) $\hat{L}^{(10)}_t$ and the median (P50) $\hat{L}^{(50)}_t$. Second, we use \textbf{conservative value estimation} via the Lower Confidence Bound (LCB). We estimate the request value using the P10 prediction:
\begin{equation}
    \hat{v}_t^{\text{LCB}} = p_{\text{in}} \cdot n_t^{\text{in}} + p_{\text{out}} \cdot \hat{L}^{(10)}_t
\end{equation}
This conservative estimate ensures we are ``90\% confident'' that the request is worth at least this much before spending quota on it. The intuition is simple: if even the pessimistic estimate of a request's value exceeds the opportunity cost threshold, we can confidently route it to the subscription.

\subsection{Unified Risk-Aware Routing Algorithm}

We combine the shadow prices and the conservative value estimate into a single unified decision rule. For each incoming request $r_t$, we calculate the net gain of routing to the subscription:

\begin{equation}
    G_t = \hat{v}_t^{\text{LCB}} - \psi(z_t) - \lambda_t \cdot w_t \cdot \hat{p}_t
\end{equation}

where $\hat{p}_t$ is the estimated duration and $w_t$ is the concurrency weight (typically 1, but larger for models requiring more GPU memory).

\begin{algorithm}[tb]
   \caption{Unified Risk-Aware Routing (LA-PD)}
   \label{alg:lapd}
\begin{algorithmic}
   \STATE {\bfseries Input:} Request $r_t$, quota usage $k_t$, system state
   \STATE {\bfseries Parameters:} Quota $Q$, value bounds $[L, U]$
   \STATE 1. Predict quantiles $\hat{L}^{(10)}_t \gets \mathcal{M}(r_t)$
   \STATE 2. Estimate conservative value $\hat{v}^{\text{LCB}}_t$
   \STATE 3. Get shadow prices $\psi(z_t)$ and $\lambda_t$
   \STATE 4. Calculate Net Gain $G_t = \hat{v}^{\text{LCB}}_t - \psi(z_t) - \lambda_t \cdot w_t \cdot \hat{p}_t$
   \IF{$G_t > 0$ \AND Capacity Available}
       \STATE Route to Subscription ($S_Q$ or $S_C$)
       \STATE Update usage $k_{t+1} \gets k_t + 1$
   \ELSE
       \STATE Route to API ($S_A$)
   \ENDIF
\end{algorithmic}
\end{algorithm}

This algorithm unifies the economic and physical constraints. When concurrency is not a bottleneck ($\lambda_t \approx 0$), it reduces to the optimal online knapsack policy. When concurrency is tight, it naturally shifts to prioritizing short, high-value jobs.
